\chapter{Plane curves}\label{geometry:ch:curves}

\begin{center}\itshape
Nine points are worth but eight.
\upshape(title of R.~E.~Bradley and L.~Stemkoski's account of Euler's letter to Cramer of 1744)
\end{center}

In 1720 the young Scottish mathematician Colin Maclaurin, in his \emph{Geometria Organica}, noticed something odd. A curve of degree $n$ in the plane is given by an equation with
$\frac{(n+1)(n+2)}2$ coefficients, defined up to a common factor, so it should be determined by $\frac{n(n+3)}2$ of its points: a line by $2$ points,
a conic by $5$, a cubic by $9$. On the other hand two curves of degrees $m$ and $n$ meet, as people already suspected, in $mn$ points. Two cubics meet in
$9$ points. But then these $9$ points cannot determine a cubic, since at least two cubics pass through them. Gabriel Cramer and Leonhard Euler,
who were writing books on curves at the time, corresponded about it, and in a letter to Cramer of 20 October 1744 Euler explained the way out: when nine
points are the intersection of two cubics, they are special, they impose only eight conditions, and infinitely many cubics pass through them. Euler
published the explanation in 1750, the year of Cramer's treatise on curves, and the puzzle is now called Cramer's paradox.

The answer turned into one of the great theorems of classical geometry, the theorem of Cayley and Bacharach: any cubic through eight of the nine
points passes through the ninth. And the count $mn$ became Bézout's theorem, published by \'Etienne B\'ezout in 1779. This chapter proves both with one
tool, the resultant, which decides by a determinant whether two polynomials have a common root. As a reward, Pascal's theorem of Appendix~D drops out
in a few lines, and the next chapter will use Cayley--Bacharach to prove that the points of a cubic curve form a group.

\begin{think}
How many points can a circle and the graph of $y=x^3-3x$ have in common? Draw them. Now try two circles: Bézout's theorem promises $2\times2=4$
intersection points, but two circles never meet in more than two. Where are the other two?
\end{think}

\section{The resultant}

Let $A$ be a commutative ring (for us a field, or a polynomial ring over a field) and
\[
f=a_0x^m+a_1x^{m-1}+\dots+a_m,\qquad g=b_0x^n+b_1x^{n-1}+\dots+b_n
\]
with coefficients in $A$, $m,n\ge1$. The \emph{Sylvester matrix} $S(f,g)$ is the $(m+n)\times(m+n)$ matrix whose rows are the coefficient vectors of
$x^{n-1}f,\dots,xf,f,x^{m-1}g,\dots,xg,g$, written in the basis $x^{m+n-1},\dots,x,1$. So it has $n$ rows of $a$'s and $m$ rows of $b$'s, each shifted one
step to the right from the row above. The \emph{resultant} is $\operatorname{Res}(f,g)=\det S(f,g)$. For instance, for $m=2$, $n=1$,
\[
\operatorname{Res}(a_0x^2+a_1x+a_2,\ b_0x+b_1)=\det\begin{pmatrix}a_0&a_1&a_2\\b_0&b_1&0\\0&b_0&b_1\end{pmatrix}.
\]

\begin{proposition}\label{geometry:prop:resultant}
\begin{enumerate}
\item There are polynomials $u,v\in A[x]$ with $\deg u<n$, $\deg v<m$ and $uf+vg=\operatorname{Res}(f,g)$.
\item If $A$ is a field and $a_0b_0\ne0$, then $\operatorname{Res}(f,g)=0$ if and only if $f$ and $g$ have a common factor of positive degree.
\end{enumerate}
\end{proposition}

\begin{proof}
A row vector
\[
c=(u_{n-1},\dots,u_0,v_{m-1},\dots,v_0)
\]
gives $u=\sum u_ix^i$, $v=\sum v_ix^i$, and $cS$ is the coefficient vector of $uf+vg$. (1) Let
$c$ be the last row of the adjugate matrix of $S$. Since $\operatorname{adj}(S)S=\det(S)I$, $cS=(0,\dots,0,\det S)$, which says
$uf+vg=\operatorname{Res}(f,g)$; the entries of the adjugate are polynomials in the coefficients.

(2) Over a field, $\det S=0$ if and only if $cS=0$ for some nonzero $c$, that is, $uf+vg=0$ with $\deg u<n$, $\deg v<m$ and $(u,v)\ne(0,0)$. If $h$ is
a common factor of positive degree, $u=g/h$ and $v=-f/h$ work. Conversely, suppose $uf=-vg$ with $(u,v)\ne0$. Then $v\ne0$, since $v=0$ would give
$uf=0$ and $u=0$. If $f$ and $g$ had no common factor, $f$ would divide $vg$ and hence $v$, impossible since $0\le\deg v<m=\deg f$.
\end{proof}

Over an algebraically closed field, a common factor of positive degree is the same as a common root. So the resultant is a single polynomial in the
coefficients that vanishes exactly when the two equations have a common solution. Exercise~\ref{geometry:ex:resproduct} gives the formula
$\operatorname{Res}(f,g)=a_0^nb_0^m\prod_{i,j}(\alpha_i-\beta_j)$ in terms of the roots.

\section{Plane curves}

Let $K$ be a field. A \emph{plane curve of degree $d$} is a nonzero homogeneous polynomial $F\in K[x,y,z]$ of degree $d$, up to a nonzero factor; its
points are the $[p]\in\mathbb P^2$ with $F(p)=0$. A line, a conic and a cubic have degrees $1$, $2$, $3$. Since $K[x,y,z]$ has unique factorisation,
$F$ is a product of irreducible factors, the \emph{components} of the curve. Homogeneous polynomials of degree $d$ form a vector space of dimension
$\binom{d+2}2=\frac{(d+1)(d+2)}2$ (the monomials $x^iy^jz^k$ with $i+j+k=d$), so curves of degree $d$ form a projective space of dimension
$\frac{d(d+3)}2$. Asking a curve to pass through a point is one linear condition on the coefficients, and therefore:

\begin{proposition}\label{geometry:prop:dimcurves}
Through any $\frac{d(d+3)}2$ points of $\mathbb P^2(K)$ passes a curve of degree $d$.
\end{proposition}

A line meets a curve of degree $d$ in at most $d$ points, unless it is a component: restricted to the line, $F$ becomes a binary form of degree $d$, which
either is zero or has at most $d$ roots.

A point $P$ of the curve $F$ is \emph{singular} if $F_x(P)=F_y(P)=F_z(P)=0$, and \emph{smooth} otherwise; at a smooth point the \emph{tangent line} is
$F_x(P)x+F_y(P)y+F_z(P)z=0$. By Euler's formula $xF_x+yF_y+zF_z=dF$, in an affine chart $z=1$ a point is singular if and only if $f=f_x=f_y=0$, where
$f(x,y)=F(x,y,1)$. Figure~\ref{geometry:fig:singular} shows the three kinds of singular cubics: a node, where two branches cross, a cusp, and the node of the folium of
Descartes.

\begin{figure}[ht]
\centering
\fig{geometry-ch07-singular}
\caption{Singular cubics: a node $y^2=x^2(x+1)$, a cusp $y^2=x^3$, and the folium of Descartes $x^3+y^3=3xy$, whose node is at the origin. The singular
points are in red.}
\label{geometry:fig:singular}
\end{figure}

\section{B\'ezout's theorem}

To intersect two curves $F$ and $G$, we eliminate one variable with the resultant. Assume $O=[0:0:1]$ is on neither curve. Then the coefficient of
$z^m$ in $F$ is $F(0,0,1)\ne0$, and we can write
\[
F=a_0z^m+a_1(x,y)z^{m-1}+\dots+a_m(x,y),\qquad G=b_0z^n+b_1(x,y)z^{n-1}+\dots+b_n(x,y),
\]
with $a_0,b_0$ nonzero constants and $a_k,b_k$ homogeneous of degree $k$ in $x,y$. Let $R(x,y)=\operatorname{Res}_z(F,G)$, the resultant with
respect to $z$, a polynomial in $x$ and $y$.

\begin{lemma}\label{geometry:lem:reshomog}
$R$ is homogeneous of degree $mn$, or zero.
\end{lemma}

\begin{proof}
Index the rows and columns of the Sylvester matrix from $0$. In row $r<n$ the entry in column $c$ is $a_{c-r}$ (zero if $c-r\notin[0,m]$), and in row
$n+r$ it is $b_{c-r}$. Replacing $(x,y)$ by $(tx,ty)$ multiplies $a_k$ and $b_k$ by $t^k$. Multiply row $r$ and row $n+r$ by $t^r$, and divide column $c$
by $t^c$: the entry $a_{c-r}t^{c-r}$ becomes $a_{c-r}$, and similarly for $b$. Hence
\[
R(tx,ty)\cdot t^{\sum_{r<n}r+\sum_{r<m}r}=t^{\sum_{c<m+n}c}\,R(x,y),
\]
and $\frac{(m+n)(m+n-1)}2-\frac{n(n-1)}2-\frac{m(m-1)}2=mn$.
\end{proof}

\begin{theorem}[B\'ezout]\label{geometry:thm:bezout}
Let $F$ and $G$ be plane curves of degrees $m$ and $n$ over a field $K$, with no common component. Then they meet in at most $mn$ points.
If $K$ is algebraically closed, then for a suitable choice of coordinates each common point $P$ has a multiplicity $I_P(F,G)\ge1$, defined below, and
\[
\sum_{P\in F\cap G}I_P(F,G)=mn.
\]
\end{theorem}

\begin{proof}
We first assume $K$ algebraically closed, and come back to the general case at the end.

\emph{The resultant is not zero.} Choose coordinates with $O=[0:0:1]$ on neither curve; this is possible because the nonzero polynomial $FG$ does not
vanish on all of $\mathbb P^2(K)$ ($K$ is infinite). If $R=0$, Proposition~\ref{geometry:prop:resultant} applied over the field $K(x,y)$ gives a common factor of
$F$ and $G$ of positive degree in $z$ in $K(x,y)[z]$, and by Gauss's lemma a common factor in $K[x,y][z]$, which is excluded. So $R$ is a nonzero form of
degree $mn$ by Lemma~\ref{geometry:lem:reshomog}.

\emph{Roots of $R$ are lines through $O$.} Fix $(x_0,y_0)\ne(0,0)$. Since $a_0,b_0$ are nonzero constants, $R(x_0,y_0)$ is the resultant of the
one-variable polynomials $F(x_0,y_0,z)$ and $G(x_0,y_0,z)$ of degrees $m$ and $n$, so it vanishes if and only if they have a common root $z_0$, that is, if and only
if the line through $O$ and $[x_0:y_0:0]$ contains a common point $[x_0:y_0:z_0]$ of $F$ and $G$.

\emph{Finiteness.} So $F\cap G$ is contained in at most $mn$ lines through $O$. Choose a second point $O'$ on neither curve and on none of these lines;
the same argument puts $F\cap G$ in at most $mn$ lines through $O'$. A line through $O$ and a line through $O'$ are distinct, so they meet in one point, and
$F\cap G$ has at most $(mn)^2$ points.

\emph{Counting.} Now choose $O$ on neither curve and on none of the finitely many lines joining two points of $F\cap G$. Then each line through $O$
contains at most one point of $F\cap G$, so $P\mapsto(\text{the line }OP)$ is a bijection from $F\cap G$ to the roots of $R$ in $\mathbb P^1$. There are at
most $mn$ of them. Define $I_P(F,G)$ as the multiplicity of the corresponding root of $R$. Over an algebraically closed field the form $R$ is a product of
$mn$ linear factors, so the multiplicities add up to $mn$.

\emph{Other fields.} Over any field $K$, the points of $F\cap G$ over $K$ are among the points over an algebraic closure $\overline K$,
and $F,G$ still have no common component over $\overline K$ (a common factor over $\overline K$ would give, through the resultant, $R=0$, and $R$ does not
depend on the field). So there are at most $mn$ of them.
\end{proof}

The multiplicity was defined using a choice of $O$. In fact it does not depend on that choice, and it can be defined intrinsically as the dimension of a
local ring (Fulton's \emph{Algebraic Curves}, Appendix~C); we will not need this. We do need to know that tangency and singular points count at least twice.
Call $O$ \emph{admissible} for $F$ and $G$ if it is on neither curve, on no line joining two common points, and on none of the tangent lines to $G$ at
common points where $G$ is smooth. There are only finitely many lines to avoid, so admissible points exist.

\begin{lemma}\label{geometry:lem:singularmult}
Let $O$ be admissible and $P\in F\cap G$. If $P$ is a singular point of $F$, then $I_P(F,G)\ge2$.
\end{lemma}

\begin{proof}
A linear change of the coordinates $x,y$ (leaving $z$ alone) fixes $O$ and changes $R$ by the same substitution, so we may assume $P=[0:1:z_0]$; then the
line $OP$ is $x=0$ and the root of $R$ corresponding to $P$ is $[0:1]$. Set $y=1$ and write $f(x,z)=F(x,1,z)$, $g(x,z)=G(x,1,z)$ and $r(x)=R(x,1)$. By
Proposition~\ref{geometry:prop:resultant}, $r(x)=u(x,z)f(x,z)+v(x,z)g(x,z)$ for some polynomials $u,v$, for all $x$ and $z$. At the point $(0,z_0)$ we have $f=g=0$, and
$f_x=f_z=0$ because $P$ is singular on $F$. Differentiating the identity with respect to $z$ (the left side does not depend on $z$) and with respect to
$x$, at $(0,z_0)$:
\[
0=v(0,z_0)\,g_z(0,z_0),\qquad r'(0)=v(0,z_0)\,g_x(0,z_0).
\]
If $g_z(0,z_0)\ne0$, the first equation gives $v(0,z_0)=0$, so $r'(0)=0$. If $g_z(0,z_0)=0$ and $g_x(0,z_0)=0$, then $r'(0)=0$ directly. The remaining case,
$g_z=0\ne g_x$, means that $G$ is smooth at $P$ with tangent line $x=0$, the line $OP$, which admissibility excludes. So $r(0)=r'(0)=0$: the root is at
least double.
\end{proof}

\begin{figure}[ht]
\centering
\fig{geometry-ch07-bezout}
\caption{The circle $x^2+y^2=3.24$ and the cubic $y=x^3-3x$ meet in $6=2\times3$ points, the most B\'ezout's theorem allows.}
\label{geometry:fig:bezout}
\end{figure}

The two circles of the opening question meet in at most $2$ points of the affine plane, and B\'ezout's theorem promises $4$. The missing two are at
infinity: every circle passes through the two \emph{circular points} $[1:\pm i:0]$ (Exercise~\ref{geometry:ex:circularpoints}). B\'ezout's count is honest only
in the complex projective plane.

\section{How many singular points?}

A curve can have many singular points if it is reducible: $d$ lines in general position have $\binom d2$ crossing points. Irreducible curves are
much more restricted.

\begin{theorem}\label{geometry:thm:singbound}
An irreducible plane curve of degree $d$ over an algebraically closed field has at most $\frac{(d-1)(d-2)}2$ singular points.
\end{theorem}

\begin{proof}
For $d=1$ there is nothing to prove, and for $d=2$ an irreducible conic is smooth (Exercise~\ref{geometry:ex:irreducibleconic}). Let $d\ge3$ and suppose $F$ has
$N=\frac{(d-1)(d-2)}2+1$ singular points. Choose also $d-3$ further points of $F$, smooth and distinct (an irreducible curve of degree $d\ge2$ has
infinitely many points: every line through a point of it meets it again, by the factorisation of the binary form). That is
$N+d-3=\frac{(d-2)(d+1)}2$ points, while curves of degree $d-2$ form a projective space of dimension $\frac{(d-2)(d+1)}2$. By
Proposition~\ref{geometry:prop:dimcurves} some curve $G$ of degree $d-2$ passes through all of them. It has no common component with $F$, which is irreducible of
larger degree. Choose an admissible $O$. By Lemma~\ref{geometry:lem:singularmult}, each singular point counts at least twice, so
\[
\sum_PI_P(F,G)\ge2N+(d-3)=(d-1)(d-2)+2+d-3=d(d-2)+1,
\]
which contradicts B\'ezout's theorem: the sum is exactly $d(d-2)$.
\end{proof}

The bound is sharp: the nodal cubic has $1=\frac{2\cdot1}2$ singular point, and the lemniscate of Bernoulli is an irreducible quartic with $3$
(Exercise~\ref{geometry:ex:lemniscate}). Curves that reach the bound can be parametrised by rational functions, like the circle in
Appendix~D and the nodal cubic in Exercise~\ref{geometry:ex:nodalparam}. More generally, $\frac{(d-1)(d-2)}2$ minus the number of singular points, each counted with a
suitable weight, is the \emph{genus} of the curve, which Volume~12 studies, and the curves of genus $0$ are exactly the ones with a rational parametrisation.

\section{Cayley--Bacharach}

We now resolve Cramer's paradox. Recall that points \emph{impose independent conditions} on curves of degree $d$ if the linear conditions ``$F(p_i)=0$''
are linearly independent, equivalently if for each $i$ there is a curve of degree $d$ through all the points except $p_i$.

\begin{lemma}\label{geometry:lem:fivepointsconics}
Five points of $\mathbb P^2(K)$, no four of them collinear, impose independent conditions on conics.
\end{lemma}

\begin{proof}
Let $p$ be one of the points and $q_1,\dots,q_4$ the others; we need a conic through the $q_i$ missing $p$. If three of the $q_i$ lie on a line $L$, say
$q_1,q_2,q_3$, then $p\notin L$ (no four collinear); take $L$ together with a line through $q_4$ missing $p$. Otherwise no three of the $q_i$ are collinear.
If $p$ is not on the pair of lines $q_1q_2\cup q_3q_4$, we are done. If it is, say on $q_1q_2$, use $q_1q_3\cup q_2q_4$ instead: $p\in q_1q_3$ would make $p$ the
intersection $q_1$ of $q_1q_2$ and $q_1q_3$, and $p\in q_2q_4$ would make $p=q_2$; both are impossible.
\end{proof}

\begin{lemma}\label{geometry:lem:eightpoints}
Eight distinct points of $\mathbb P^2(K)$, no four collinear and no seven on a conic, impose independent conditions on cubics.
\end{lemma}

\begin{proof}
By symmetry it suffices to find a cubic through $p_1,\dots,p_7$ that misses $p_8$.

\emph{Case A: three of $p_1,\dots,p_7$ lie on a line $L$.} Then $p_8\notin L$. The other four of $p_1,\dots,p_7$, together with $p_8$, are five points with no four
collinear, so by Lemma~\ref{geometry:lem:fivepointsconics} there is a conic $Q$ through the four that misses $p_8$. The cubic $L\cup Q$ works.

\emph{Case B: no three of $p_1,\dots,p_7$ are collinear.} Then any five of them lie on a unique conic, which is nondegenerate (Theorem~\ref{geometry:thm:fivepoints}).
If six of them lie on a conic $Q$, then $p_8\notin Q$ (no seven on a conic), and $Q$ together with a line through the seventh point that misses $p_8$ works. So
assume no six of $p_1,\dots,p_7$ are on a conic. For a pair $\{i,j\}\subset\{1,\dots,7\}$ let $S$ be the complementary five points and $C_S$ the conic through
them; the cubic $C_S\cup p_ip_j$ passes through $p_1,\dots,p_7$. If it misses $p_8$ for some pair, we are done. The pairs whose line contains $p_8$ are
pairwise disjoint (two of them sharing $p_i$ would put three of $p_1,\dots,p_7$ on the line $p_8p_i$), so there are at most $3$ of them, and at least $18$ of
the $21$ pairs have $p_8\notin p_ip_j$. Suppose that for all of these $p_8\in C_S$. Among $18$ pairs, two share a point (pairwise disjoint pairs from
seven points number at most $3$), and then their complements $S,S'$ share four points. The conics $C_S$ and $C_{S'}$ both contain these four points and $p_8$;
two distinct nondegenerate conics meet in at most $4$ points by B\'ezout, so $C_S=C_{S'}$, which then contains $S\cup S'$, six of the points
$p_1,\dots,p_7$. This was excluded.
\end{proof}

\begin{theorem}[Cayley--Bacharach for cubics]\label{geometry:thm:cayleybacharach}
Let $C_1$ and $C_2$ be cubics over a field with no common component, meeting in nine distinct points. Then every cubic through eight of the
nine points passes through the ninth.
\end{theorem}

\begin{proof}
No four of the nine points are on a line $L$: otherwise $L$ would meet $C_1$ in four points, so $L$ would be a component of $C_1$, and likewise of $C_2$.
No seven are on a conic $Q$: if $Q$ is irreducible, B\'ezout ($7>6$) makes it a component of $C_1$ and of $C_2$; if $Q$ is a pair of lines, one of them
contains four of the points. So any eight of the points satisfy Lemma~\ref{geometry:lem:eightpoints}, and the cubics through them form a vector space of dimension
$10-8=2$. It contains the forms $C_1$ and $C_2$, which are not proportional, so it is exactly the pencil $\lambda C_1+\mu C_2$. Every member vanishes at the ninth
point.
\end{proof}

This is the resolution of Cramer's paradox: the nine common points of two cubics impose only eight conditions on cubics, so they lie on a whole pencil of
them. Here is the reward.

\begin{corollary}[Pascal, again]\label{geometry:cor:pascalcb}
Let $A,B,C,D,E,F$ be six points of a nondegenerate conic $Q$. Then $X=AB\cap DE$, $Y=BC\cap EF$ and $Z=CD\cap FA$ are collinear.
\end{corollary}

\begin{proof}
Let $C_1=AB\cup CD\cup EF$ and $C_2=BC\cup DE\cup FA$, two cubics made of three lines each (Figure~\ref{geometry:fig:pascalcb}). Each line of $C_1$ meets each line of
$C_2$ once, and the nine meeting points are $A,B,C,D,E,F,X,Y,Z$. They are distinct (a line meets $Q$ in at most two points, so no three of the six are
collinear, and none of $X,Y,Z$ is one of the six), and the cubics share no line. The cubic $Q\cup XY$ passes through the eight points $A,\dots,F,X,Y$, so by
Theorem~\ref{geometry:thm:cayleybacharach} it passes through $Z$. But $Z$ is on $CD$, which meets $Q$ only at $C$ and $D$, so $Z\notin Q$. Hence $Z\in XY$.
\end{proof}

\begin{figure}[ht]
\centering
\fig{geometry-ch07-pascalcb}
\caption{Pascal by Cayley--Bacharach. The blue cubic $AB\cup CD\cup EF$ and the red cubic $BC\cup DE\cup FA$ meet in nine points. The conic together with the
line $XY$ is a cubic through eight of them, so it contains the ninth, $Z$.}
\label{geometry:fig:pascalcb}
\end{figure}

\begin{remark}
The general Cayley--Bacharach theorem replaces cubics by curves of degrees $m$ and $n$ meeting in $mn$ points: a curve of degree $m+n-3$ through all but one of
them passes through the last. It follows from Max Noether's ``$AF+BG$'' theorem, proved in Volume~12, Chapter~11. The history, from Chasles and Cayley to
Bacharach, and the modern view are told by Eisenbud, Green and Harris in the article listed below. In Appendix~G, Theorem~\ref{geometry:thm:cayleybacharach} gives the
associativity of the group law on a cubic.
\end{remark}

\section{Exercises}

\begin{exercise}[$\star$]\label{geometry:ex:reslinear}
Show that $\operatorname{Res}(ax^2+bx+c,\ x-t)=at^2+bt+c$.
\end{exercise}

\begin{exercise}[$\star$]\label{geometry:ex:discriminant}
Compute $\operatorname{Res}(f,f')$ for $f=x^2+px+q$, and explain why it vanishes exactly when $f$ has a double root.
\end{exercise}

\begin{exercise}[$\star\star$]\label{geometry:ex:resproduct}
Let $f=a_0\prod_{i=1}^m(x-\alpha_i)$ and $g=b_0\prod_{j=1}^n(x-\beta_j)$. Prove that $\operatorname{Res}(f,g)=a_0^nb_0^m\prod_{i,j}(\alpha_i-\beta_j)$.
\end{exercise}

\begin{exercise}[$\star\star$]\label{geometry:ex:circularpoints}
Show that every circle $x^2+y^2+2axz+2byz+cz^2=0$ passes through the circular points $[1:\pm i:0]$. Find the four intersections promised by B\'ezout
for two non-concentric circles that meet in two real points, and for two concentric circles.
\end{exercise}

\begin{exercise}[$\star$]\label{geometry:ex:folium}
Show that the folium $x^3+y^3=3xyz$ has exactly one singular point in $\mathbb P^2(\mathbb C)$, find its two tangent lines, and parametrise the
curve by the lines $y=tx$ through it.
\end{exercise}

\begin{exercise}[$\star$]\label{geometry:ex:nodalparam}
Parametrise the nodal cubic $y^2=x^2(x+1)$ and the cuspidal cubic $y^2=x^3$ by rational functions, using the lines through the singular point.
\end{exercise}

\begin{exercise}[$\star\star$]\label{geometry:ex:lemniscate}
Find the singular points of the lemniscate $(x^2+y^2)^2=x^2-y^2$ in $\mathbb P^2(\mathbb C)$ and compare with Theorem~\ref{geometry:thm:singbound}.
\end{exercise}

\begin{exercise}[$\star$]\label{geometry:ex:irreducibleconic}
Show that a conic with a singular point over an algebraically closed field is a pair of lines (possibly equal). Deduce that irreducible conics are smooth.
\end{exercise}

\begin{exercise}[$\star\star$]\label{geometry:ex:polarcurve}
Let $F$ be a smooth curve of degree $d\ge2$ over $\mathbb C$ and $Q=[q]$ a point. Show that the points of $F$ whose tangent passes through $Q$ lie on the
\emph{polar curve} $q_0F_x+q_1F_y+q_2F_z=0$, and deduce that at most $d(d-1)$ tangents to $F$ pass through a point not on $F$: six for a smooth cubic.
\end{exercise}

\begin{exercise}[$\star$]\label{geometry:ex:cramergrid}
The cubics $x^3-xz^2$ and $y^3-yz^2$ meet in the nine points $\{-1,0,1\}^2$ of the affine plane. Find all cubics through these nine points.
\end{exercise}

\begin{exercise}[$\star\star$]\label{geometry:ex:cbpappus}
Deduce Pappus' theorem (Theorem~\ref{geometry:thm:pappusline}) from Theorem~\ref{geometry:thm:cayleybacharach}.
\end{exercise}

\begin{exercise}[$\star\star\star$]\label{geometry:ex:braikenridge}
(Braikenridge and Maclaurin) Let $A,\dots,F$ be six points, no three collinear, such that the nine points $A,\dots,F$, $X=AB\cap DE$, $Y=BC\cap EF$,
$Z=CD\cap FA$ are distinct. Show that if $X,Y,Z$ are collinear, the six points lie on a conic: the converse of Pascal's theorem.
\end{exercise}

\section{Hints}

\begin{hint}{geometry:ex:reslinear}
Expand the $3\times3$ Sylvester determinant along the first row.
\end{hint}
\begin{hint}{geometry:ex:discriminant}
$f'=2x+p$; the Sylvester matrix is $3\times3$.
\end{hint}
\begin{hint}{geometry:ex:resproduct}
Treat the $\alpha_i,\beta_j$ as indeterminates. The resultant vanishes when $\alpha_i=\beta_j$; compare degrees, then evaluate at $\alpha=0$.
\end{hint}
\begin{hint}{geometry:ex:circularpoints}
Subtract the equations of two circles: what is left factors, one factor being $z$.
\end{hint}
\begin{hint}{geometry:ex:folium}
Compute the three partial derivatives. The lowest-degree part of the equation at the origin gives the tangent lines.
\end{hint}
\begin{hint}{geometry:ex:nodalparam}
Substitute $y=tx$ and cancel the double root $x=0$.
\end{hint}
\begin{hint}{geometry:ex:lemniscate}
Homogenise, then look at $z=0$ and at $z=1$ separately.
\end{hint}
\begin{hint}{geometry:ex:irreducibleconic}
At a singular point $P$ of $q(x)=x^{\mathsf T}Ax$, $AP=0$. Move $P$ to $[0:0:1]$.
\end{hint}
\begin{hint}{geometry:ex:polarcurve}
The tangent at $P$ is $\sum F_i(P)x_i=0$. For the count, use B\'ezout; a smooth curve is irreducible.
\end{hint}
\begin{hint}{geometry:ex:cramergrid}
Check the hypotheses of Lemma~\ref{geometry:lem:eightpoints} for eight of the nine points.
\end{hint}
\begin{hint}{geometry:ex:cbpappus}
Use the hexagon $AB'CA'BC'$ of Exercise~\ref{geometry:ex:pascalpappus}, and the cubic $\ell\cup\ell'\cup XY$.
\end{hint}
\begin{hint}{geometry:ex:braikenridge}
Take the conic through $A,B,C,D,E$ and the Pascal line: together they are a cubic through eight of the nine points.
\end{hint}

\section{Solutions}

\begin{solution}{geometry:ex:reslinear}
\idea The Sylvester matrix is small enough to expand.
\writeup With $m=2$, $n=1$ the rows are $(a,b,c)$, $(1,-t,0)$, $(0,1,-t)$, and
$\det=a(t^2-0)-b(-t-0)+c(1-0)=at^2+bt+c$. This is $f(t)$, as it should be: $x-t$ and $f$ have a common root exactly when $f(t)=0$.
\end{solution}

\begin{solution}{geometry:ex:discriminant}
\idea A double root is a common root of $f$ and $f'$.
\writeup The rows are $(1,p,q)$, $(2,p,0)$, $(0,2,p)$, and $\det=1\cdot p^2-p\cdot2p+q\cdot4=4q-p^2$. This vanishes exactly when the discriminant
$p^2-4q$ does. By Proposition~\ref{geometry:prop:resultant}, it vanishes if and only if $f$ and $f'$ have a common root, which happens exactly when $f$ has a double root.
\end{solution}

\begin{solution}{geometry:ex:resproduct}
\idea A polynomial that vanishes whenever $\alpha_i=\beta_j$ is divisible by $\alpha_i-\beta_j$; the degrees then leave room for only a constant.
\writeup Work in the polynomial ring over $\mathbb Z$ in the variables $a_0$, $b_0$, $\alpha_1,\dots,\alpha_m$ and $\beta_1,\dots,\beta_n$, a unique factorisation domain. The coefficients are
$a_k=(-1)^ka_0e_k(\alpha)$ and $b_k=(-1)^kb_0e_k(\beta)$, where $e_k$ are the elementary symmetric polynomials. The resultant $\rho$ is a polynomial in these
variables. Setting $\alpha_i=\beta_j$ gives two polynomials with a common root, so $\rho$ vanishes; as a polynomial in $\alpha_i$ it is therefore divisible by
$\alpha_i-\beta_j$. These $mn$ factors are distinct primes, so their product $\pi$ divides $\rho$. The argument of Lemma~\ref{geometry:lem:reshomog} shows that every term
of the determinant has total weight $mn$ when $a_k,b_k$ have weight $k$; since $a_k$ has degree $k$ in the $\alpha$'s, $\rho$ has total degree at most $mn$ in the
$\alpha$'s and $\beta$'s, like $\pi$. So $\rho=c\,\pi$ with $c\in\mathbb Z[a_0,b_0]$. Each term of the determinant contains $n$ entries from the $a$-rows and $m$ from the
$b$-rows, so $c=\lambda a_0^nb_0^m$ for an integer $\lambda$. Finally set all $\alpha_i=0$: then $f=a_0x^m$, and the Sylvester matrix is block triangular with
$a_0$ on the first $n$ diagonal entries and a lower triangular block with $b_n$ on the diagonal, so $\rho=a_0^nb_n^m$. Since $b_n=(-1)^nb_0\prod\beta_j$, this is
$a_0^nb_0^m\prod_j(-\beta_j)^m=a_0^nb_0^m\prod_{i,j}(0-\beta_j)$, so $\lambda=1$.
\end{solution}

\begin{solution}{geometry:ex:circularpoints}
\idea All circles share two points at infinity, which is why two circles have only two ``real'' intersections left.
\writeup At $z=0$ the equation becomes $x^2+y^2=0$, whose solutions are $[1:i:0]$ and $[1:-i:0]$. For two circles $C$ and $C'$, the difference of the equations is
$z\,\ell$ with $\ell=2(a-a')x+2(b-b')y+(c-c')z$. So $C\cap C'=C\cap(\{z=0\}\cup\{\ell=0\})$: the two circular points, plus the points of $C$ on the line $\ell=0$, the
\emph{radical axis}. If the circles meet in two real points, these are the two points on the radical axis, and with the circular points we get $2+2=4$. For
concentric circles $a=a'$ and $b=b'$, so the difference is $(c-c')z^2$: the circles meet only at the circular points, each counted twice. Their tangents there
coincide (the circles are tangent at infinity), which is why the multiplicity is $2$.
\end{solution}

\begin{solution}{geometry:ex:folium}
\idea Singular points are common zeros of the three partial derivatives.
\writeup With $F=x^3+y^3-3xyz$: $F_x=3x^2-3yz$, $F_y=3y^2-3xz$, $F_z=-3xy$. From $F_z=0$, $x=0$ or $y=0$. If $x=0$ then $F_y=3y^2=0$, so $y=0$; if $y=0$ then
$x=0$. The only candidate is $[0:0:1]$, which is on the curve. In the chart $z=1$ the equation is $-3xy+(x^3+y^3)=0$; its lowest-degree part $-3xy$ shows two
branches with tangents $x=0$ and $y=0$: a node. Substituting $y=tx$ gives $x^3(1+t^3)=3tx^2$, so apart from the double root $x=0$ the line meets the curve at
$x=\frac{3t}{1+t^3}$, $y=\frac{3t^2}{1+t^3}$.
\end{solution}

\begin{solution}{geometry:ex:nodalparam}
\idea A line through a double point meets a cubic in only one more point.
\writeup For $y^2=x^2(x+1)$, $y=tx$ gives $t^2x^2=x^2(x+1)$, so $x=t^2-1$ and $y=t(t^2-1)$. For $y^2=x^3$, $y=tx$ gives $t^2x^2=x^3$, so $x=t^2$, $y=t^3$. In both cases
every point except the singular one corresponds to exactly one $t$.
\end{solution}

\begin{solution}{geometry:ex:lemniscate}
\idea The lemniscate has a node at the origin and two more singular points hiding at the circular points.
\writeup Homogenise: $F=(x^2+y^2)^2-z^2(x^2-y^2)$, with $F_x=4x(x^2+y^2)-2xz^2$, $F_y=4y(x^2+y^2)+2yz^2$, $F_z=-2z(x^2-y^2)$. At $z=0$, $F_z=0$ and $F_x,F_y$
vanish when $x^2+y^2=0$: the circular points $[1:\pm i:0]$, which lie on the curve. At $z=1$, $F_z=0$ gives $x^2=y^2$. If $y\ne0$, $F_y=0$ gives
$x^2+y^2=-\frac12$, so $x^2=-\frac14\ne0$, and then $F_x=x(4(x^2+y^2)-2)=-4x\ne0$; so $y=0$, then $x=0$: the origin, a node. That is three singular points, the
maximum $\frac{3\cdot2}2$ allowed for an irreducible quartic.

The lemniscate is indeed irreducible. One checks that $x=\frac{t+t^3}{1+t^4}$, $y=\frac{t-t^3}{1+t^4}$ satisfies the equation: $x^2+y^2=\frac{2t^2}{1+t^4}$ and
$x^2-y^2=(x+y)(x-y)=\frac{4t^4}{(1+t^4)^2}$. Since $K(t)$ has no zero divisors, some irreducible factor $G$ of $F$ vanishes on this parametrisation. The line
$x=\frac12$ meets it where $t^4-2t^3-2t+1=0$; dividing by $t^2$ and setting $u=t+\frac1t$ gives $u^2-2u-2=0$, so $u=1\pm\sqrt3$, and the four roots $t$ are
distinct. They give four distinct points, because $t\mapsto(x,y)$ is injective away from $t=0$: the ratio $\frac yx=\frac{1-t^2}{1+t^2}$ determines $t^2$, and
$x$ changes sign with $t$. So the line $x=\frac12$, which is not a component, meets $G$ in four points, $\deg G\ge4$, and $G=F$.
\end{solution}

\begin{solution}{geometry:ex:irreducibleconic}
\idea A singular quadratic form has a kernel, and a binary quadratic form splits over an algebraically closed field.
\writeup For $F=x^{\mathsf T}Ax$, the gradient is $2Ax$, so $P=[p]$ is singular exactly when $Ap=0$. Choose coordinates with $p=e_3$. Then the third row and column of
$A$ are zero, and $F$ is a binary quadratic form in $x,y$, which factors into two linear forms over an algebraically closed field: two lines through $P$, or a
double line. So a conic with a singular point is reducible, and an irreducible conic is smooth.
\end{solution}

\begin{solution}{geometry:ex:polarcurve}
\idea Tangency through $Q$ is one more equation, of degree $d-1$.
\writeup The tangent at a smooth point $P$ is $\sum_iF_i(P)x_i=0$, so it passes through $Q$ if and only if $\sum_iq_iF_i(P)=0$: $P$ lies on the polar curve
$G=q_0F_x+q_1F_y+q_2F_z$, of degree $d-1$. The form $G$ is not zero for a suitable $Q$: if $F_x,F_y,F_z$ were linearly dependent, then after a linear change of
coordinates $F$ would not involve one variable, and it would be a union of lines through a point, which is singular for $d\ge2$. A smooth curve is irreducible
(two components would meet, by B\'ezout, and a common point of two components is singular), and $\deg G<\deg F$, so $F$ and $G$ have no common component. By
B\'ezout they meet in at most $d(d-1)$ points, so at most $d(d-1)$ tangents pass through $Q$. For a smooth cubic, at most $6$; one finds numerically
exactly $6$ for random $Q$.
\end{solution}

\begin{solution}{geometry:ex:cramergrid}
\idea This is Cramer's paradox in miniature: nine points, one pencil.
\writeup The two cubics are $x(x-z)(x+z)$ and $y(y-z)(y+z)$, three vertical and three horizontal lines with no common component, meeting in the nine grid
points. A line contains at most three of them, so no four are collinear. A conic meets each horizontal line $y=c$ in at most two points unless it contains it; a
conic containing seven grid points would therefore contain a horizontal line and then lie on a pair of lines, which contain at most six grid points. So
any eight of the nine points impose independent conditions by Lemma~\ref{geometry:lem:eightpoints}, and the cubics through all nine form exactly the pencil
$\lambda(x^3-xz^2)+\mu(y^3-yz^2)$, as in the proof of Theorem~\ref{geometry:thm:cayleybacharach}.
\end{solution}

\begin{solution}{geometry:ex:cbpappus}
\idea Pappus is Pascal for the degenerate conic $\ell\cup\ell'$, and Cayley--Bacharach does not mind degenerate conics.
\writeup With $A,B,C\in\ell$ and $A',B',C'\in\ell'$ (all different from $\ell\cap\ell'$), take the cubics $C_1=AB'\cup CA'\cup BC'$ and $C_2=B'C\cup A'B\cup C'A$.
Intersecting each line of $C_1$ with each line of $C_2$ gives the nine points $B',X,A$, then $C,A',Y$, then $Z,B,C'$, where $X=AB'\cap A'B$, $Y=AC'\cap A'C$,
$Z=BC'\cap B'C$. Assume these nine points are distinct (the general case; the degenerate cases follow by a limiting argument or by Appendix~C). The cubic
$\ell\cup\ell'\cup XY$ passes through $A,B,C,A',B',C',X,Y$, so it contains $Z$. If $Z\in\ell$, then $Z$, being on $BC'$, would be $B$, and then $B$ would lie on
$B'C$, a line meeting $\ell$ only at $C$; so $Z\notin\ell$, and similarly $Z\notin\ell'$. Hence $Z\in XY$.
\end{solution}

\begin{solution}{geometry:ex:braikenridge}
\idea Run the proof of Corollary~\ref{geometry:cor:pascalcb} backwards: the conic through five of the points and the Pascal line make a cubic through eight of the nine points.
\writeup Let $C_1=AB\cup CD\cup EF$ and $C_2=BC\cup DE\cup FA$. As in Corollary~\ref{geometry:cor:pascalcb} they meet exactly in the nine distinct points
$A,\dots,F,X,Y,Z$, and they share no line (no three of the six points are collinear). Let $Q$ be the conic through $A,B,C,D,E$ and $\ell$ the line containing
$X,Y,Z$. The cubic $Q\cup\ell$ passes through the eight points $A,B,C,D,E,X,Y,Z$, so by Theorem~\ref{geometry:thm:cayleybacharach} it passes through $F$. If $F$ were on
$\ell$, then $\ell$ would contain $F$ and $Z$, two distinct points of the line $FA$, so $\ell=FA\ni A$; and it would contain $F$ and $Y$, so $\ell=EF\ni E$; then
$A,E,F$ would be collinear. So $F\in Q$: the six points lie on a conic.
\end{solution}

\section*{Sources and further reading}

\begin{itemize}
\item R.~E.~Bradley and L.~Stemkoski, \emph{When nine points are worth but eight: Euler's resolution of Cramer's paradox}, Convergence (MAA), online (Maclaurin's
\emph{Geometria Organica} of 1720, Euler's letter to Cramer of 20 October 1744, and Euler's paper \emph{Sur une contradiction apparente dans la doctrine des
lignes courbes}, M\'em.\ Acad.\ Berlin 4, 1750; the epigraph of this chapter shortens the title of Bradley and Stemkoski's article); G.~Cramer, \emph{Introduction \`a l'analyse des lignes courbes
alg\'ebriques}, Geneva, 1750.
\item \'E.~B\'ezout, \emph{Th\'eorie g\'en\'erale des \'equations alg\'ebriques}, Paris, 1779.
\item D.~Eisenbud, M.~Green and J.~Harris, \emph{Cayley--Bacharach theorems and conjectures}, Bulletin of the AMS 33 (1996), 295--324 (the history and the proof
of Lemma~\ref{geometry:lem:eightpoints} in this spirit).
\item W.~Fulton, \emph{Algebraic Curves}, Benjamin, 1969, freely available from the author (intersection multiplicity, B\'ezout, Noether's theorem).
\item F.~Kirwan, \emph{Complex Algebraic Curves}, Cambridge University Press, 1992; E.~Brieskorn and H.~Kn\"orrer, \emph{Plane Algebraic Curves}, Birkh\"auser,
1986 (a big, beautiful and historical book).
\item C.~G.~Gibson, \emph{Elementary Geometry of Algebraic Curves}, Cambridge University Press, 1998 (resultants and singular points, with many pictures).
\end{itemize}
